\chapter{Abstract}

\subsection*{Part~\ref{part:neutron-stars} (project thesis): Neutron Stars}

Neutron stars contain matter at extremely high densities, where both general
relativity and quantum many-body physics are needed to describe their structure.
In this project, we study how different models of dense matter affect the
mass--radius relation of neutron stars. The analysis is based on the
Tolman--Oppenheimer--Volkoff equations for hydrostatic equilibrium, which are
solved numerically for different equations of state.
We first consider an ideal, non-interacting Fermi gas. This illustrates how
degeneracy pressure can counteract a gravitational collapse, while also showing
the limitations of such idealized models when compared with observed
neutron-star properties. We then develop an interacting model for nucleons and
leptons in charge neutrality and \(\beta\)-equilibrium, where the interactions
are described by self-consistent meson fields. This gives a more realistic
equation of state, which is combined with a standard crust model before being
used in the stellar calculations.
The resulting mass--radius relations give radii typical of neutron stars, but
do not support the heaviest observed stars and differ from current astronomical
constraints. Thus, while the interacting model improves substantially on the
ideal Fermi-gas description, it remains incomplete at the highest densities. The
project therefore motivates the need for more advanced equations of state to
describe the full range of observed neutron-star properties.

\subsection*{Part~\ref{part:quark-stars} (Master's thesis): Quark Stars}

Compact-star solutions are constructed from two-flavor effective quark-matter
models, with emphasis on the Quark--Meson--Diquark model and two-flavor color
superconductivity. The equations of state are obtained in the mean-field
approximation and used as input to the Tolman--Oppenheimer--Volkoff equations.
The two-flavor Quark--Meson model is first used as a reference case. Varying
the sigma mass changes both the stiffness of the equation of state and the
shape of the stellar branch: lower sigma masses give larger maximum masses and
self-bound-like low-mass behavior, while the \(m_\sigma=600~\mathrm{MeV}\)
case gives a gravitationally bound sequence. A Bodmer--Witten-inspired vacuum
shift reduces the mass and radius scale, but does not qualitatively change this
ordering in the explored range.
The main result is obtained from the Quark--Meson--Diquark model with 2SC
pairing. Retaining the full one-loop potential is essential for obtaining a
robust paired phase. For neutral stellar matter, the default QMD parameter set
gives a gravitationally bound sequence with \(M_{\max}=1.970\,M_\odot\) and
\(R(M_{\max})=12.16\,\mathrm{km}\).
Parameter variations show that \(g_\Delta\), \(m_\Delta\), and \(\lambda_3\)
control complementary aspects of the mass--radius relation. Selected combined
parameter choices exceed \(2\,M_\odot\) and are compatible with the plotted
NICER regions for both J0030+0451 and J0740+6620. The results indicate that
two-flavor QMD matter with 2SC pairing can produce phenomenologically promising
compact-star sequences, but the present models are pure quark-star sequences
without a hadronic envelope, strange quarks, or tidal deformabilities.